\chapter{Special Cases: Multiple Optima, Degeneracy, Unboundedness,
  Infeasibility}

For a well-behaved maximisation problem the simplex method of the previous
chapters marches from corner to corner and stops in a final tableau that
describes exactly one optimal vertex. This chapter treats the four ways an
LP can deviate from that script: the objective may be maximised along a
whole edge instead of a single corner (multiple optima); a corner may be pinned down by more constraints
than necessary (degeneracy); the feasible region may stretch
to infinity in a direction along which the objective keeps growing
(unboundedness); or the constraints may contradict
one another outright (infeasibility).

Each case leaves an unmistakable fingerprint in \emph{two} places: in the
graphical picture obtained by corner enumeration, and in the simplex
tableau. Reading these fingerprints quickly (and not confusing them with
one another) is an exam skill in its own right.

\section{Core Rules}

\begin{rulebox}{Tableau conventions (recap of Chapters 4--5)}
Maximisation tableaus are written with one row per basic variable and the
objective row $z$ at the bottom. The RHS column holds the values of the
basic variables and, in the $z$-row, the current value of $z$.
\begin{itemize}[leftmargin=1.6em,itemsep=0.3ex]
  \item \textbf{Optimality:} the tableau is optimal as soon as every
        $z$-row entry is $\ge 0$.
  \item \textbf{Entering variable:} a column with a negative $z$-row entry
        (most negative first). Increasing that variable by one unit raises
        $z$ by the absolute value of its entry.
  \item \textbf{Leaving variable:} ratio test
        $\min_i \{\text{RHS}_i/a_i : a_i>0\}$ over the \emph{positive}
        entries $a_i$ of the entering column only.
  \item We call the $z$-row entry of a \emph{nonbasic} column its
        \emph{reduced cost}; basic columns always carry the entry $0$.
\end{itemize}
\end{rulebox}

\begin{rulebox}{Multiple optima}
\textbf{In the tableau:} the tableau is optimal, but some \emph{nonbasic}
variable has reduced cost $0$. That variable may enter without changing
$z$: after the usual ratio test and pivot, the new tableau shows a
\emph{different} optimal vertex with the \emph{same} objective value. If
$x^{A}$ and $x^{B}$ are the two optimal vertices, the entire connecting
edge is optimal, and the \emph{general optimal solution} is
\[
  x(\lambda)\;=\;\lambda\,x^{A}+(1-\lambda)\,x^{B},
  \qquad 0\le\lambda\le 1 .
\]
\textbf{Graphically:} the level lines of $z$ are parallel to a binding edge
of the feasible region; corner enumeration returns the maximal value at two
\emph{adjacent} corners, and every point between them is optimal too.
(In more than two variables the optimal set can be a whole face: the set
of all convex combinations of the optimal vertices; in the plane it is
always a segment.)
\end{rulebox}

\begin{rulebox}{Degeneracy}
A basic feasible solution is \emph{degenerate} if at least one \emph{basic}
variable equals $0$ (a zero in the RHS column above the $z$-row).
\begin{itemize}[leftmargin=1.6em,itemsep=0.3ex]
  \item \textbf{In the tableau:} degeneracy is created by a \emph{tie in
        the ratio test}. The leaving row is not unique; whichever row is
        chosen, the basic variable of the other tying row stays basic at
        value $0$ in the next tableau.
  \item \textbf{Geometrically:} at the corresponding vertex \emph{more}
        constraints are active than the dimension requires, in $\R^{2}$
        three or more boundary lines pass through one point. The vertex is
        \emph{overdetermined}, and several different bases describe the same
        point.
  \item \textbf{Consequences:} a pivot may have step length $0$ (the basis
        changes, the point and $z$ do not). In artificial examples the
        simplex method can even \emph{cycle} forever among the bases of one
        vertex; anti-cycling rules (e.g.\ Bland's rule: among the eligible
        candidates always pick the smallest index) exclude this. In hand
        computations degeneracy is harmless, simply carry on.
  \item \textbf{Do not confuse:} degeneracy $=$ a \emph{basic} variable at
        zero; multiple optima $=$ a \emph{nonbasic} reduced cost at zero.
        Neither implies the other.
\end{itemize}
\end{rulebox}

\begin{rulebox}{Unboundedness}
\textbf{In the tableau:} an entering column (negative reduced cost)
contains \emph{no positive entry}, so no ratio can be formed. The entering
variable can then be increased without limit while feasibility is
preserved, and $z\to\infty$: the LP is unbounded, it has no finite
optimum. The improving ray can be read off directly: set the entering
variable to $t\ge 0$; each basic variable changes by $-\,(\text{column
entry})\cdot t$.\\[0.4ex]
\textbf{Graphically:} the feasible region is unbounded \emph{and} $z$
improves along the escape direction. Beware: an unbounded region alone does
not make the LP unbounded: the objective must keep growing along it.
\end{rulebox}

\begin{rulebox}{Infeasibility}
Detected in \textbf{phase one}: minimise the sum of the artificial
variables, $w=\sum a_i$ (equivalently, maximise $-w$ with the usual
maximisation rules; the bottom row then shows $-w$).
\begin{itemize}[leftmargin=1.6em,itemsep=0.3ex]
  \item If $\min w=0$: a feasible basis has been found; drop the
        artificials and start phase two.
  \item If $\min w>0$ (some artificial variable is \emph{stuck basic at
        a positive value} in the optimal phase-one tableau) then no
        feasible solution exists: the feasible region is empty. The original objective never
        plays a role.
\end{itemize}
\textbf{Graphically:} the constraint half-planes have no common point.
\end{rulebox}

\begin{rulebox}{Diagnosis at a glance}
\begin{center}
\begin{tabular}{p{2.9cm}|p{5.7cm}|p{5.6cm}}
\hline
\textbf{Case} & \textbf{In the tableau} & \textbf{In the picture}\\
\hline
unique optimum & optimal; all nonbasic reduced costs $>0$ &
  the $z$-line touches the region in exactly one corner\\
\hline
multiple optima & optimal; a \emph{nonbasic} reduced cost equals $0$ &
  $z$-line parallel to a binding edge; two adjacent corners tie for the
  maximum\\
\hline
degenerate vertex & a \emph{basic} variable equals $0$ (tie in the ratio
  test one step earlier) & more lines through the corner than needed\\
\hline
unbounded LP & entering column without positive entries &
  region open in a direction along which $z$ grows\\
\hline
infeasible LP & phase one ends with $w>0$ (artificial stuck positive) &
  the constraints have no common point\\
\hline
\end{tabular}
\end{center}
\end{rulebox}

\section{Practice Problems}

All variables are nonnegative. Tableaus follow the conventions of the recap
box: objective row at the bottom, and a maximisation tableau is optimal
when all entries of that row are nonnegative.

\begin{problems}
  \item Solve graphically by corner enumeration:
        \[
          \max\; z=2x_1+2x_2
          \quad\text{s.t.}\quad
          x_1+x_2\le 4,\qquad x_1\le 3,\qquad x_2\le 3 .
        \]
        List all corners with their objective values, explain why the LP
        has multiple optima, write down the general optimal solution, and
        verify that the midpoint of the optimal edge is optimal as well.

  \item Solve by the simplex method:
        \[
          \max\; z=2x_1+x_2
          \quad\text{s.t.}\quad
          2x_1+x_2\le 8,\qquad x_1+x_2\le 5 .
        \]
        Show that the final tableau signals multiple optima, pivot once
        more to reach a second optimal vertex, and give the set of all
        optimal solutions.

  \item For the LP
        $\max\; z=x_1+2x_2$ s.t.\ $x_1+2x_2\le 6$, $x_1+x_2\le 4$,
        the simplex method has produced the tableau
        \begin{center}
        \begin{tabular}{c|rrrr|r}
        \hline
        Basis & $x_1$ & $x_2$ & $s_1$ & $s_2$ & RHS\\
        \hline
        $x_2$ & $\tfrac{1}{2}$ & $1$ & $\tfrac{1}{2}$ & $0$ & $3$\\
        $s_2$ & $\tfrac{1}{2}$ & $0$ & $-\tfrac{1}{2}$ & $1$ & $1$\\
        \hline
        $z$ & $0$ & $0$ & $1$ & $0$ & $6$\\
        \hline
        \end{tabular}
        \end{center}
        Diagnose the tableau, read off the current solution, and determine
        \emph{all} optimal solutions of the LP.

  \item Solve by the simplex method:
        \[
          \max\; z=3x_1+x_2
          \quad\text{s.t.}\quad
          2x_1+x_2\le 4,\qquad x_1\le 2,\qquad x_2\le 3 .
        \]
        Point out the tie in the ratio test, show that the next tableau is
        degenerate, interpret the degeneracy geometrically, and decide
        whether the optimum is unique.

  \item The following tableau arose in a maximisation problem:
        \begin{center}
        \begin{tabular}{c|rrrrr|r}
        \hline
        Basis & $x_1$ & $x_2$ & $s_1$ & $s_2$ & $s_3$ & RHS\\
        \hline
        $x_1$ & $1$ & $0$ & $1$ & $-1$ & $0$ & $2$\\
        $x_2$ & $0$ & $1$ & $-1$ & $2$ & $0$ & $0$\\
        $s_3$ & $0$ & $0$ & $2$ & $-1$ & $1$ & $3$\\
        \hline
        $z$ & $0$ & $0$ & $1$ & $1$ & $0$ & $8$\\
        \hline
        \end{tabular}
        \end{center}
        (a) Is the tableau optimal? (b) Is the basic solution degenerate,
        and if so, which variable shows it? (c) Read off $x_1$, $x_2$ and
        $z$. (d) Is the optimum unique? (e) What does the degeneracy mean
        geometrically?

  \item Consider
        \[
          \max\; z=2x_1+x_2
          \quad\text{s.t.}\quad
          x_1-x_2\le 1,\qquad -x_1+x_2\le 2 .
        \]
        (a) Describe the feasible region and argue graphically that the LP
        is unbounded. (b) Run the simplex method until the unboundedness
        criterion fires. (c) Give a feasible ray along which $z\to\infty$.

  \item During a simplex computation (maximisation) the following tableau
        appeared:
        \begin{center}
        \begin{tabular}{c|rrrr|r}
        \hline
        Basis & $x_1$ & $x_2$ & $s_1$ & $s_2$ & RHS\\
        \hline
        $x_1$ & $1$ & $-2$ & $1$ & $0$ & $3$\\
        $s_2$ & $0$ & $-1$ & $2$ & $1$ & $4$\\
        \hline
        $z$ & $0$ & $-4$ & $3$ & $0$ & $12$\\
        \hline
        \end{tabular}
        \end{center}
        Diagnose the tableau, exhibit the improving ray, and show that $z$
        is unbounded along it.

  \item Consider
        \[
          \max\; z=3x_1+2x_2
          \quad\text{s.t.}\quad
          x_1+x_2\le 2,\qquad x_1+x_2\ge 4 .
        \]
        Explain why phase one is needed, carry it out, and state the
        conclusion. What does the picture of the two constraints show?

  \item Run phase one for
        \[
          \max\; z=x_1+x_2
          \quad\text{s.t.}\quad
          x_1+x_2\le 2,\qquad x_1+2x_2\ge 6 ,
        \]
        interpret the final value of $w$, and explain the result
        graphically.

  \item For which values of $c\ge 0$ does the LP
        \[
          \max\; z=c\,x_1+x_2
          \quad\text{s.t.}\quad
          x_1+x_2\le 4,\qquad 2x_1+x_2\le 6
        \]
        have multiple optimal solutions? For each such $c$ give the set of
        all optima.

  \item Let $0\le c\le 5$. For which $c$ does
        \[
          \max\; z=4x_1+c\,x_2
          \quad\text{s.t.}\quad
          x_1+x_2\le 6,\qquad x_1\le 4
        \]
        have more than one optimal solution? Describe the optimal set as a
        function of $c$ on the whole interval $[0,5]$.

  \item Classify each of the following maximisation tableaus: optimal or
        not, unique or multiple optima, degenerate or not, or unbounded.
        Read off the solution where one exists.
        \begin{enumerate}[label=(\alph*),leftmargin=1.8em,itemsep=0.5ex]
          \item \begin{tabular}[t]{c|rrrr|r}
                \hline
                Basis & $x_1$ & $x_2$ & $s_1$ & $s_2$ & RHS\\
                \hline
                $x_1$ & $1$ & $0$ & $2$ & $-1$ & $3$\\
                $x_2$ & $0$ & $1$ & $-1$ & $1$ & $2$\\
                \hline
                $z$ & $0$ & $0$ & $1$ & $2$ & $12$\\
                \hline
                \end{tabular}
          \item \begin{tabular}[t]{c|rrrr|r}
                \hline
                Basis & $x_1$ & $x_2$ & $s_1$ & $s_2$ & RHS\\
                \hline
                $x_1$ & $1$ & $2$ & $1$ & $0$ & $4$\\
                $s_2$ & $0$ & $1$ & $-1$ & $1$ & $0$\\
                \hline
                $z$ & $0$ & $3$ & $2$ & $0$ & $8$\\
                \hline
                \end{tabular}
          \item \begin{tabular}[t]{c|rrrr|r}
                \hline
                Basis & $x_1$ & $x_2$ & $s_1$ & $s_2$ & RHS\\
                \hline
                $s_1$ & $-2$ & $0$ & $1$ & $1$ & $3$\\
                $x_2$ & $-1$ & $1$ & $0$ & $2$ & $4$\\
                \hline
                $z$ & $-3$ & $0$ & $0$ & $1$ & $8$\\
                \hline
                \end{tabular}
        \end{enumerate}

  \item Solve completely with the simplex method and determine \emph{all}
        optimal solutions:
        \[
          \max\; z=3x_1+3x_2
          \quad\text{s.t.}\quad
          x_1+x_2\le 5,\qquad 2x_1+x_2\le 8,\qquad x_2\le 4 .
        \]
        Check the result against a corner enumeration.

  \item Decide without lengthy computation whether each LP is infeasible,
        unbounded, or has a finite optimum. Justify briefly; for (a) also
        argue via the phase-one objective.
        \begin{enumerate}[label=(\alph*),leftmargin=1.8em,itemsep=0.3ex]
          \item $\max\; x_1+x_2$ \;s.t.\; $x_1-x_2\ge 1$,\; $x_2-x_1\ge 1$
          \item $\max\; x_1+x_2$ \;s.t.\; $-x_1+x_2\le 1$
        \end{enumerate}
\end{problems}

\section{Solutions}

\begin{solutions}
  \item The feasible region has the corners $(0,0)$, $(3,0)$, $(3,1)$,
        $(1,3)$, $(0,3)$. Corner enumeration:
        \begin{center}
        \begin{tabular}{c|ccccc}
        \hline
        corner & $(0,0)$ & $(3,0)$ & $(3,1)$ & $(1,3)$ & $(0,3)$\\
        \hline
        $z=2x_1+2x_2$ & $0$ & $6$ & $8$ & $8$ & $6$\\
        \hline
        \end{tabular}
        \end{center}
        The maximum $z^{*}=8$ is attained at \emph{two adjacent} corners,
        $(3,1)$ and $(1,3)$. Reason: the level lines $2x_1+2x_2=\text{const}$
        have slope $-1$, exactly the slope of the binding edge
        $x_1+x_2=4$: objective and edge are parallel. General optimal
        solution:
        \[
          (x_1,x_2)=\lambda\,(3,1)+(1-\lambda)\,(1,3)
                   =(1+2\lambda,\;3-2\lambda),\qquad \lambda\in[0,1],
        \]
        all with $z^{*}=8$. Midpoint ($\lambda=\tfrac{1}{2}$): $(2,2)$ is
        feasible ($2+2=4\le 4$, $2\le 3$) and $z=2\cdot 2+2\cdot 2=8$.
        \emph{Optimal.}

  \item Standard form: $2x_1+x_2+s_1=8$, $x_1+x_2+s_2=5$. Start tableau:
        \begin{center}
        \begin{tabular}{c|rrrr|r}
        \hline
        Basis & $x_1$ & $x_2$ & $s_1$ & $s_2$ & RHS\\
        \hline
        $s_1$ & $2$ & $1$ & $1$ & $0$ & $8$\\
        $s_2$ & $1$ & $1$ & $0$ & $1$ & $5$\\
        \hline
        $z$ & $-2$ & $-1$ & $0$ & $0$ & $0$\\
        \hline
        \end{tabular}
        \end{center}
        Enter $x_1$ ($-2$); ratios $8/2=4$, $5/1=5$; $s_1$ leaves; pivot
        on the $2$:
        \begin{center}
        \begin{tabular}{c|rrrr|r}
        \hline
        Basis & $x_1$ & $x_2$ & $s_1$ & $s_2$ & RHS\\
        \hline
        $x_1$ & $1$ & $\tfrac{1}{2}$ & $\tfrac{1}{2}$ & $0$ & $4$\\
        $s_2$ & $0$ & $\tfrac{1}{2}$ & $-\tfrac{1}{2}$ & $1$ & $1$\\
        \hline
        $z$ & $0$ & $0$ & $1$ & $0$ & $8$\\
        \hline
        \end{tabular}
        \end{center}
        All $z$-row entries are $\ge 0$: optimal with $x^{A}=(4,0)$,
        $z^{*}=8$. But the \emph{nonbasic} variable $x_2$ has reduced cost
        $0$, multiple optima. Enter $x_2$; ratios
        $4/\tfrac{1}{2}=8$, $1/\tfrac{1}{2}=2$; $s_2$ leaves:
        \begin{center}
        \begin{tabular}{c|rrrr|r}
        \hline
        Basis & $x_1$ & $x_2$ & $s_1$ & $s_2$ & RHS\\
        \hline
        $x_1$ & $1$ & $0$ & $1$ & $-1$ & $3$\\
        $x_2$ & $0$ & $1$ & $-1$ & $2$ & $2$\\
        \hline
        $z$ & $0$ & $0$ & $1$ & $0$ & $8$\\
        \hline
        \end{tabular}
        \end{center}
        Second optimal vertex $x^{B}=(3,2)$, again $z^{*}=8$
        ($2\cdot 3+2=8$). All optimal solutions:
        \[
          (x_1,x_2)=\lambda\,(4,0)+(1-\lambda)\,(3,2)
                   =(3+\lambda,\;2-2\lambda),\qquad \lambda\in[0,1].
        \]

  \item The tableau is optimal (all $z$-row entries $\ge 0$); it shows
        $x_2=3$, $s_2=1$, hence the vertex $(x_1,x_2)=(0,3)$ with
        $z^{*}=6$. The nonbasic variable $x_1$ has reduced cost $0$, so the
        optimum is \emph{not} unique. Enter $x_1$; ratios
        $3/\tfrac{1}{2}=6$, $1/\tfrac{1}{2}=2$; $s_2$ leaves; pivot on the
        $\tfrac{1}{2}$ in the $s_2$-row:
        \begin{center}
        \begin{tabular}{c|rrrr|r}
        \hline
        Basis & $x_1$ & $x_2$ & $s_1$ & $s_2$ & RHS\\
        \hline
        $x_2$ & $0$ & $1$ & $1$ & $-1$ & $2$\\
        $x_1$ & $1$ & $0$ & $-1$ & $2$ & $2$\\
        \hline
        $z$ & $0$ & $0$ & $1$ & $0$ & $6$\\
        \hline
        \end{tabular}
        \end{center}
        Second optimal vertex $(2,2)$, $z=2+2\cdot 2=6$. All optima:
        \[
          (x_1,x_2)=\lambda\,(0,3)+(1-\lambda)\,(2,2)
                   =(2-2\lambda,\;2+\lambda),\qquad \lambda\in[0,1],
        \]
        with $z^{*}=6$. (Graphically: $z=x_1+2x_2$ is parallel to the
        binding constraint $x_1+2x_2=6$.)

  \item Start tableau ($s_1,s_2,s_3$ slacks of the three constraints):
        \begin{center}
        \begin{tabular}{c|rrrrr|r}
        \hline
        Basis & $x_1$ & $x_2$ & $s_1$ & $s_2$ & $s_3$ & RHS\\
        \hline
        $s_1$ & $2$ & $1$ & $1$ & $0$ & $0$ & $4$\\
        $s_2$ & $1$ & $0$ & $0$ & $1$ & $0$ & $2$\\
        $s_3$ & $0$ & $1$ & $0$ & $0$ & $1$ & $3$\\
        \hline
        $z$ & $-3$ & $-1$ & $0$ & $0$ & $0$ & $0$\\
        \hline
        \end{tabular}
        \end{center}
        Enter $x_1$ ($-3$). Ratios: $4/2=2$, $2/1=2$, no ratio in row
        $s_3$: a \emph{tie} between rows $s_1$ and $s_2$: the next
        basic solution will be degenerate. Choose the $s_1$-row (smaller
        index); pivot on the $2$:
        \begin{center}
        \begin{tabular}{c|rrrrr|r}
        \hline
        Basis & $x_1$ & $x_2$ & $s_1$ & $s_2$ & $s_3$ & RHS\\
        \hline
        $x_1$ & $1$ & $\tfrac{1}{2}$ & $\tfrac{1}{2}$ & $0$ & $0$ & $2$\\
        $s_2$ & $0$ & $-\tfrac{1}{2}$ & $-\tfrac{1}{2}$ & $1$ & $0$ & $0$\\
        $s_3$ & $0$ & $1$ & $0$ & $0$ & $1$ & $3$\\
        \hline
        $z$ & $0$ & $\tfrac{1}{2}$ & $\tfrac{3}{2}$ & $0$ & $0$ & $6$\\
        \hline
        \end{tabular}
        \end{center}
        The tableau is optimal: $x^{*}=(2,0)$, $z^{*}=6$. The basic
        variable $s_2=0$ makes the solution \emph{degenerate}.
        Geometrically, \emph{three} lines pass through the optimal corner
        $(2,0)$: $2x_1+x_2=4$, $x_1=2$ and $x_2=0$: in the plane two
        would suffice, so the vertex is overdetermined. The optimum is
        \emph{unique}: the nonbasic reduced costs $\tfrac{1}{2}$ (for
        $x_2$) and $\tfrac{3}{2}$ (for $s_1$) are strictly positive.

        \emph{Remark.} Pivoting instead in the $s_2$-row gives
        \begin{center}
        \begin{tabular}{c|rrrrr|r}
        \hline
        Basis & $x_1$ & $x_2$ & $s_1$ & $s_2$ & $s_3$ & RHS\\
        \hline
        $x_1$ & $1$ & $0$ & $0$ & $1$ & $0$ & $2$\\
        $s_1$ & $0$ & $1$ & $1$ & $-2$ & $0$ & $0$\\
        $s_3$ & $0$ & $1$ & $0$ & $0$ & $1$ & $3$\\
        \hline
        $z$ & $0$ & $-1$ & $0$ & $3$ & $0$ & $6$\\
        \hline
        \end{tabular}
        \end{center}
        which is \emph{not yet} optimal: $x_2$ must enter, the ratio test
        gives $\min\{0/1,\,3/1\}=0$, and the pivot has step length $0$,
        the basis changes but the point $(2,0)$ and $z=6$ do not. One such
        zero-step pivot later the optimality test is passed. This is the
        typical (harmless) cost of degeneracy.

  \item (a) Yes: every $z$-row entry is $\ge 0$.
        (b) Yes: the \emph{basic} variable $x_2$ has value $0$.
        (c) $x_1=2$, $x_2=0$, $z=8$ (and $s_3=3$).
        (d) Unique: the nonbasic variables $s_1$ and $s_2$ have reduced
        costs $1$ and $1$, both strictly positive; a zero would be needed
        for multiple optima. Degeneracy (a \emph{basic} zero) is a
        different phenomenon and does not create alternative optima.
        (e) At the optimal vertex more constraints are active than the
        dimension requires (here additionally $x_2=0$); several bases
        describe this one overdetermined corner.

  \item (a) The constraints say $x_1-1\le x_2\le x_1+2$: the region is the
        strip between two parallel lines of slope $1$, cut to the first
        quadrant, with corners $(0,0)$, $(1,0)$, $(0,2)$. It is unbounded
        in the direction $d=(1,1)$: from any feasible point, moving along
        $d$ keeps both constraints satisfied while
        $z=2x_1+x_2$ grows by $3$ per step: no corner is optimal.

        (b) Start tableau:
        \begin{center}
        \begin{tabular}{c|rrrr|r}
        \hline
        Basis & $x_1$ & $x_2$ & $s_1$ & $s_2$ & RHS\\
        \hline
        $s_1$ & $1$ & $-1$ & $1$ & $0$ & $1$\\
        $s_2$ & $-1$ & $1$ & $0$ & $1$ & $2$\\
        \hline
        $z$ & $-2$ & $-1$ & $0$ & $0$ & $0$\\
        \hline
        \end{tabular}
        \end{center}
        Enter $x_1$ ($-2$); only row $s_1$ has a positive entry, ratio
        $1/1=1$; pivot there:
        \begin{center}
        \begin{tabular}{c|rrrr|r}
        \hline
        Basis & $x_1$ & $x_2$ & $s_1$ & $s_2$ & RHS\\
        \hline
        $x_1$ & $1$ & $-1$ & $1$ & $0$ & $1$\\
        $s_2$ & $0$ & $0$ & $1$ & $1$ & $3$\\
        \hline
        $z$ & $0$ & $-3$ & $2$ & $0$ & $2$\\
        \hline
        \end{tabular}
        \end{center}
        Now $x_2$ must enter (reduced cost $-3$), but its column contains
        the entries $-1$ and $0$, \emph{no positive entry}. The
        unboundedness criterion fires: the LP has no finite optimum.

        (c) Set $x_2=t\ge 0$: the tableau gives $x_1=1+t$, $s_2=3$, so
        \[
          \begin{pmatrix}x_1\\ x_2\end{pmatrix}
          =\begin{pmatrix}1\\ 0\end{pmatrix}
           +t\begin{pmatrix}1\\ 1\end{pmatrix},
          \qquad z=2+3t\;\longrightarrow\;\infty .
        \]
        (Check: $(1+t)-t=1$, $-(1+t)+t=-1\le 2$, feasible for every
        $t\ge 0$.)

  \item The tableau is not optimal: $x_2$ has reduced cost $-4$ and must
        enter. Its column contains $-2$ and $-1$: no positive entry, so
        the ratio test fails: the LP is \emph{unbounded}. Improving ray:
        set $x_2=t\ge 0$; then $x_1=3+2t$ and $s_2=4+t$ stay nonnegative
        for all $t$, and
        \[
          z=12+4t\;\longrightarrow\;\infty .
        \]
        Geometrically the feasible region has an unbounded edge in the
        direction $(2,1)$ of the $(x_1,x_2)$-plane, and the objective grows
        without bound along it.

  \item The origin violates $x_1+x_2\ge 4$, so there is no all-slack
        starting basis: subtract a surplus variable $s_2$ and add an
        artificial variable $a$ in the second row, then run phase one,
        $\min w=a$ (we maximise $-w$; after expressing $a$ through the
        nonbasic variables, $-w=-4+x_1+x_2-s_2$). Start tableau:
        \begin{center}
        \begin{tabular}{c|rrrrr|r}
        \hline
        Basis & $x_1$ & $x_2$ & $s_1$ & $s_2$ & $a$ & RHS\\
        \hline
        $s_1$ & $1$ & $1$ & $1$ & $0$ & $0$ & $2$\\
        $a$   & $1$ & $1$ & $0$ & $-1$ & $1$ & $4$\\
        \hline
        $-w$ & $-1$ & $-1$ & $0$ & $1$ & $0$ & $-4$\\
        \hline
        \end{tabular}
        \end{center}
        Enter $x_1$ (smallest index among the tied entries $-1,-1$);
        ratios $2/1=2$, $4/1=4$; $s_1$ leaves:
        \begin{center}
        \begin{tabular}{c|rrrrr|r}
        \hline
        Basis & $x_1$ & $x_2$ & $s_1$ & $s_2$ & $a$ & RHS\\
        \hline
        $x_1$ & $1$ & $1$ & $1$ & $0$ & $0$ & $2$\\
        $a$   & $0$ & $0$ & $-1$ & $-1$ & $1$ & $2$\\
        \hline
        $-w$ & $0$ & $0$ & $1$ & $1$ & $0$ & $-2$\\
        \hline
        \end{tabular}
        \end{center}
        All entries of the $-w$-row are $\ge 0$: phase one is finished with
        $\min w=2>0$. The artificial variable is stuck basic at $a=2$:
        the LP is \emph{infeasible}, and phase two (with the objective
        $3x_1+2x_2$) never starts. The picture confirms it: the half-plane
        below the line $x_1+x_2=2$ and the half-plane above the parallel
        line $x_1+x_2=4$ have no common point.

  \item Standard form: $x_1+x_2+s_1=2$ and $x_1+2x_2-s_2+a=6$ with
        artificial $a$; phase one minimises $w=a$, i.e.\
        $-w=-6+x_1+2x_2-s_2$:
        \begin{center}
        \begin{tabular}{c|rrrrr|r}
        \hline
        Basis & $x_1$ & $x_2$ & $s_1$ & $s_2$ & $a$ & RHS\\
        \hline
        $s_1$ & $1$ & $1$ & $1$ & $0$ & $0$ & $2$\\
        $a$   & $1$ & $2$ & $0$ & $-1$ & $1$ & $6$\\
        \hline
        $-w$ & $-1$ & $-2$ & $0$ & $1$ & $0$ & $-6$\\
        \hline
        \end{tabular}
        \end{center}
        Enter $x_2$ ($-2$); ratios $2/1=2$, $6/2=3$; $s_1$ leaves:
        \begin{center}
        \begin{tabular}{c|rrrrr|r}
        \hline
        Basis & $x_1$ & $x_2$ & $s_1$ & $s_2$ & $a$ & RHS\\
        \hline
        $x_2$ & $1$ & $1$ & $1$ & $0$ & $0$ & $2$\\
        $a$   & $-1$ & $0$ & $-2$ & $-1$ & $1$ & $2$\\
        \hline
        $-w$ & $1$ & $0$ & $2$ & $1$ & $0$ & $-2$\\
        \hline
        \end{tabular}
        \end{center}
        Optimal for phase one, with $\min w=2>0$ and $a=2$ still basic:
        \emph{infeasible}. Interpretation: within $x_1+x_2\le 2$,
        $x\ge 0$, the largest achievable value of $x_1+2x_2$ is $4$ (at
        the point $(0,2)$); the requirement $x_1+2x_2\ge 6$ is missed by
        exactly $w^{*}=2$. Graphically the two constraint sets do not
        intersect.

  \item Corners of the feasible region: $(0,0)$, $(3,0)$, $(2,2)$ (from
        $x_1+x_2=4$ and $2x_1+x_2=6$) and $(0,4)$. Values of
        $z=c\,x_1+x_2$:
        \begin{center}
        \begin{tabular}{c|cccc}
        \hline
        corner & $(0,0)$ & $(3,0)$ & $(2,2)$ & $(0,4)$\\
        \hline
        $z$ & $0$ & $3c$ & $2c+2$ & $4$\\
        \hline
        \end{tabular}
        \end{center}
        Multiple optima require a tie between two \emph{adjacent} corners
        that also attains the maximum. The level lines of $z$ have slope
        $-c$; the edges have slopes $-2$ (on $2x_1+x_2=6$) and $-1$ (on
        $x_1+x_2=4$).
        \begin{itemize}[leftmargin=1.6em,itemsep=0.3ex]
          \item $c=1$: $2c+2=4$, tie between $(2,2)$ and $(0,4)$ with
                $z^{*}=4$ (and $(3,0)$ gives only $3$). Optimal set:
                $(x_1,x_2)=\lambda(2,2)+(1-\lambda)(0,4)
                =(2\lambda,\,4-2\lambda)$, $\lambda\in[0,1]$.
          \item $c=2$: $3c=2c+2=6$, tie between $(3,0)$ and $(2,2)$ with
                $z^{*}=6$ (and $(0,4)$ gives only $4$). Optimal set:
                $(x_1,x_2)=\lambda(3,0)+(1-\lambda)(2,2)
                =(2+\lambda,\,2-2\lambda)$, $\lambda\in[0,1]$.
        \end{itemize}
        For all other $c\ge 0$ the optimum is unique ($0\le c<1$:
        $(0,4)$; $1<c<2$: $(2,2)$; $c>2$: $(3,0)$). Answer:
        $c\in\{1,2\}$.

  \item Corners: $(0,0)$, $(4,0)$, $(4,2)$, $(0,6)$, with
        \[
          z(4,0)=16,\qquad z(4,2)=16+2c,\qquad z(0,6)=6c .
        \]
        For $0<c<4$ we have $16+2c>16$ and $16+2c>6c$, so $(4,2)$ is the
        unique optimum. Ties:
        \begin{itemize}[leftmargin=1.6em,itemsep=0.3ex]
          \item $c=0$: $(4,0)$ and $(4,2)$ both give $z^{*}=16$; the
                optimal set is the vertical edge
                $\{(4,t):0\le t\le 2\}$ on $x_1=4$.
          \item $c=4$: $(4,2)$ and $(0,6)$ both give $z^{*}=24$; the
                optimal set is
                $(x_1,x_2)=\lambda(4,2)+(1-\lambda)(0,6)
                =(4\lambda,\,6-4\lambda)$, $\lambda\in[0,1]$, on the edge
                $x_1+x_2=6$.
          \item $4<c\le 5$: unique optimum $(0,6)$, $z^{*}=6c$.
        \end{itemize}
        Multiple optima occur exactly for $c=0$ and $c=4$.

  \item
        \begin{enumerate}[label=(\alph*),leftmargin=1.8em,itemsep=0.3ex]
          \item Optimal (all $z$-entries $\ge 0$), nondegenerate (basic
                values $3,2>0$), and \emph{unique}: the nonbasic reduced
                costs $1$ and $2$ are strictly positive. Solution
                $x=(3,2)$, $z^{*}=12$.
          \item Optimal and \emph{degenerate}: the basic variable $s_2=0$.
                Still unique, because the nonbasic reduced costs $3$ (for
                $x_2$) and $2$ (for $s_1$) are positive. Solution
                $x=(4,0)$, $z^{*}=8$.
          \item Not optimal: $x_1$ has reduced cost $-3$. Its column
                contains $-2$ and $-1$, no positive entry,
                \emph{unbounded}. Along $x_1=t$: $s_1=3+2t$, $x_2=4+t$,
                $z=8+3t\to\infty$.
        \end{enumerate}

  \item Start tableau:
        \begin{center}
        \begin{tabular}{c|rrrrr|r}
        \hline
        Basis & $x_1$ & $x_2$ & $s_1$ & $s_2$ & $s_3$ & RHS\\
        \hline
        $s_1$ & $1$ & $1$ & $1$ & $0$ & $0$ & $5$\\
        $s_2$ & $2$ & $1$ & $0$ & $1$ & $0$ & $8$\\
        $s_3$ & $0$ & $1$ & $0$ & $0$ & $1$ & $4$\\
        \hline
        $z$ & $-3$ & $-3$ & $0$ & $0$ & $0$ & $0$\\
        \hline
        \end{tabular}
        \end{center}
        The entries $-3,-3$ tie; take $x_1$ (smallest index). Ratios:
        $5/1=5$, $8/2=4$, none in row $s_3$; $s_2$ leaves; pivot on the
        $2$:
        \begin{center}
        \begin{tabular}{c|rrrrr|r}
        \hline
        Basis & $x_1$ & $x_2$ & $s_1$ & $s_2$ & $s_3$ & RHS\\
        \hline
        $s_1$ & $0$ & $\tfrac{1}{2}$ & $1$ & $-\tfrac{1}{2}$ & $0$ & $1$\\
        $x_1$ & $1$ & $\tfrac{1}{2}$ & $0$ & $\tfrac{1}{2}$ & $0$ & $4$\\
        $s_3$ & $0$ & $1$ & $0$ & $0$ & $1$ & $4$\\
        \hline
        $z$ & $0$ & $-\tfrac{3}{2}$ & $0$ & $\tfrac{3}{2}$ & $0$ & $12$\\
        \hline
        \end{tabular}
        \end{center}
        Enter $x_2$ ($-\tfrac{3}{2}$). Ratios: $1/\tfrac{1}{2}=2$,
        $4/\tfrac{1}{2}=8$, $4/1=4$; $s_1$ leaves; pivot on the
        $\tfrac{1}{2}$:
        \begin{center}
        \begin{tabular}{c|rrrrr|r}
        \hline
        Basis & $x_1$ & $x_2$ & $s_1$ & $s_2$ & $s_3$ & RHS\\
        \hline
        $x_2$ & $0$ & $1$ & $2$ & $-1$ & $0$ & $2$\\
        $x_1$ & $1$ & $0$ & $-1$ & $1$ & $0$ & $3$\\
        $s_3$ & $0$ & $0$ & $-2$ & $1$ & $1$ & $2$\\
        \hline
        $z$ & $0$ & $0$ & $3$ & $0$ & $0$ & $15$\\
        \hline
        \end{tabular}
        \end{center}
        Optimal: $x^{A}=(3,2)$, $z^{*}=15$. The nonbasic variable $s_2$
        has reduced cost $0$, multiple optima. Enter $s_2$; ratios
        (positive entries only): $3/1=3$ in the $x_1$-row, $2/1=2$ in the
        $s_3$-row; $s_3$ leaves:
        \begin{center}
        \begin{tabular}{c|rrrrr|r}
        \hline
        Basis & $x_1$ & $x_2$ & $s_1$ & $s_2$ & $s_3$ & RHS\\
        \hline
        $x_2$ & $0$ & $1$ & $0$ & $0$ & $1$ & $4$\\
        $x_1$ & $1$ & $0$ & $1$ & $0$ & $-1$ & $1$\\
        $s_2$ & $0$ & $0$ & $-2$ & $1$ & $1$ & $2$\\
        \hline
        $z$ & $0$ & $0$ & $3$ & $0$ & $0$ & $15$\\
        \hline
        \end{tabular}
        \end{center}
        Second optimal vertex $x^{B}=(1,4)$, again $z=3+12=15$. All
        optimal solutions:
        \[
          (x_1,x_2)=\lambda\,(3,2)+(1-\lambda)\,(1,4)
                   =(1+2\lambda,\;4-2\lambda),\qquad\lambda\in[0,1].
        \]
        Corner enumeration confirms this: $(0,0)\mapsto 0$,
        $(4,0)\mapsto 12$, $(3,2)\mapsto 15$, $(1,4)\mapsto 15$,
        $(0,4)\mapsto 12$: the maximum sits on the edge
        $x_1+x_2=5$, to which $z=3x_1+3x_2$ is parallel.

  \item
        \begin{enumerate}[label=(\alph*),leftmargin=1.8em,itemsep=0.5ex]
          \item \emph{Infeasible.} Adding the two constraints gives
                $(x_1-x_2)+(x_2-x_1)=0\ge 2$, a contradiction. Phase-one
                view: with artificials $a_1,a_2$ in the two rows,
                $a_1=1-x_1+x_2+s_1$ and $a_2=1+x_1-x_2+s_2$, hence
                \[
                  w=a_1+a_2=2+s_1+s_2\;\ge\;2>0
                \]
                for \emph{every} choice of the variables: phase one can
                never drive $w$ to $0$ (indeed its start tableau is
                already optimal with $w=2$).
          \item \emph{Unbounded.} The region is nonempty ($(0,0)$ is
                feasible) and contains the ray $(t,0)$, $t\ge 0$, along
                which $z=t\to\infty$. In the tableau: entering $x_1$
                (reduced cost $-1$) finds only the entry $-1$ in its
                column, no positive entry, unboundedness at the very
                first step.
        \end{enumerate}
\end{solutions}

\section{Additional Problems}

A second pass over all four special cases.

\begin{problems}
  \item Solve graphically by corner enumeration and give all optimal
        solutions:
        \[
          \max\; z=4x_1+2x_2
          \quad\text{s.t.}\quad
          2x_1+x_2\le 6,\qquad x_1+x_2\le 4 .
        \]

  \item Solve by the simplex method and determine all optimal solutions:
        \[
          \max\; z=2x_1+4x_2
          \quad\text{s.t.}\quad
          x_1+2x_2\le 8,\qquad x_1+x_2\le 6 .
        \]

  \item Solve by the simplex method:
        \[
          \max\; z=2x_1+3x_2
          \quad\text{s.t.}\quad
          x_1+x_2\le 3,\qquad x_2\le 3 .
        \]
        Point out the tie in the ratio test, identify the degenerate basic
        variable in the final tableau, and interpret the situation
        geometrically. Is the optimum unique?

  \item Show that the following LP is unbounded, both from the tableau and
        by exhibiting a feasible ray with $z\to\infty$:
        \[
          \max\; z=x_1+2x_2
          \quad\text{s.t.}\quad
          -2x_1+x_2\le 1,\qquad x_1-x_2\le 2 .
        \]

  \item Run phase one for
        \[
          \max\; z=x_1+x_2
          \quad\text{s.t.}\quad
          2x_1+x_2\ge 8,\qquad x_1+x_2\le 3
        \]
        and interpret the result, including the meaning of the final value
        of $w$.

  \item For which values of $c\ge 0$ does
        \[
          \max\; z=x_1+c\,x_2
          \quad\text{s.t.}\quad
          x_1+2x_2\le 4,\qquad x_1+x_2\le 3
        \]
        have multiple optimal solutions? Give the optimal set for each such
        $c$.
\end{problems}

\section{Additional Solutions}

\begin{solutions}
  \item Corners and values of $z=4x_1+2x_2$: $(0,0)\mapsto 0$,
        $(3,0)\mapsto 12$, $(2,2)\mapsto 12$ (intersection of
        $2x_1+x_2=6$ and $x_1+x_2=4$), $(0,4)\mapsto 8$. The maximum
        $z^{*}=12$ is attained at the two adjacent corners $(3,0)$ and
        $(2,2)$: the level lines of $z$ (slope $-2$) are parallel to the
        binding edge $2x_1+x_2=6$. All optimal solutions:
        \[
          (x_1,x_2)=\lambda\,(3,0)+(1-\lambda)\,(2,2)
                   =(2+\lambda,\;2-2\lambda),\qquad \lambda\in[0,1].
        \]

  \item Start tableau:
        \begin{center}
        \begin{tabular}{c|rrrr|r}
        \hline
        Basis & $x_1$ & $x_2$ & $s_1$ & $s_2$ & RHS\\
        \hline
        $s_1$ & $1$ & $2$ & $1$ & $0$ & $8$\\
        $s_2$ & $1$ & $1$ & $0$ & $1$ & $6$\\
        \hline
        $z$ & $-2$ & $-4$ & $0$ & $0$ & $0$\\
        \hline
        \end{tabular}
        \end{center}
        Enter $x_2$ ($-4$); ratios $8/2=4$, $6/1=6$; $s_1$ leaves; pivot
        on the $2$:
        \begin{center}
        \begin{tabular}{c|rrrr|r}
        \hline
        Basis & $x_1$ & $x_2$ & $s_1$ & $s_2$ & RHS\\
        \hline
        $x_2$ & $\tfrac{1}{2}$ & $1$ & $\tfrac{1}{2}$ & $0$ & $4$\\
        $s_2$ & $\tfrac{1}{2}$ & $0$ & $-\tfrac{1}{2}$ & $1$ & $2$\\
        \hline
        $z$ & $0$ & $0$ & $2$ & $0$ & $16$\\
        \hline
        \end{tabular}
        \end{center}
        Optimal with $x^{A}=(0,4)$, $z^{*}=16$; the nonbasic $x_1$ has
        reduced cost $0$, so a second optimal vertex exists. Enter $x_1$;
        ratios $4/\tfrac{1}{2}=8$, $2/\tfrac{1}{2}=4$; $s_2$ leaves:
        \begin{center}
        \begin{tabular}{c|rrrr|r}
        \hline
        Basis & $x_1$ & $x_2$ & $s_1$ & $s_2$ & RHS\\
        \hline
        $x_2$ & $0$ & $1$ & $1$ & $-1$ & $2$\\
        $x_1$ & $1$ & $0$ & $-1$ & $2$ & $4$\\
        \hline
        $z$ & $0$ & $0$ & $2$ & $0$ & $16$\\
        \hline
        \end{tabular}
        \end{center}
        Second vertex $x^{B}=(4,2)$, $z=8+8=16$. All optimal solutions:
        \[
          (x_1,x_2)=\lambda\,(0,4)+(1-\lambda)\,(4,2)
                   =(4-4\lambda,\;2+2\lambda),\qquad \lambda\in[0,1].
        \]
        (The objective $2x_1+4x_2$ is parallel to the constraint
        $x_1+2x_2=8$.)

  \item Start tableau:
        \begin{center}
        \begin{tabular}{c|rrrr|r}
        \hline
        Basis & $x_1$ & $x_2$ & $s_1$ & $s_2$ & RHS\\
        \hline
        $s_1$ & $1$ & $1$ & $1$ & $0$ & $3$\\
        $s_2$ & $0$ & $1$ & $0$ & $1$ & $3$\\
        \hline
        $z$ & $-2$ & $-3$ & $0$ & $0$ & $0$\\
        \hline
        \end{tabular}
        \end{center}
        Enter $x_2$ ($-3$); ratios $3/1=3$ and $3/1=3$: a tie, so the
        next tableau will be degenerate. Choose the $s_1$-row:
        \begin{center}
        \begin{tabular}{c|rrrr|r}
        \hline
        Basis & $x_1$ & $x_2$ & $s_1$ & $s_2$ & RHS\\
        \hline
        $x_2$ & $1$ & $1$ & $1$ & $0$ & $3$\\
        $s_2$ & $-1$ & $0$ & $-1$ & $1$ & $0$\\
        \hline
        $z$ & $1$ & $0$ & $3$ & $0$ & $9$\\
        \hline
        \end{tabular}
        \end{center}
        Optimal: $x^{*}=(0,3)$, $z^{*}=9$, with the \emph{basic} variable
        $s_2=0$, a degenerate optimum. Geometrically the three lines
        $x_1+x_2=3$, $x_2=3$ and $x_1=0$ all pass through $(0,3)$: an
        overdetermined corner. The optimum is unique, since the nonbasic
        reduced costs $1$ (for $x_1$) and $3$ (for $s_1$) are strictly
        positive.

  \item Start tableau:
        \begin{center}
        \begin{tabular}{c|rrrr|r}
        \hline
        Basis & $x_1$ & $x_2$ & $s_1$ & $s_2$ & RHS\\
        \hline
        $s_1$ & $-2$ & $1$ & $1$ & $0$ & $1$\\
        $s_2$ & $1$ & $-1$ & $0$ & $1$ & $2$\\
        \hline
        $z$ & $-1$ & $-2$ & $0$ & $0$ & $0$\\
        \hline
        \end{tabular}
        \end{center}
        Enter $x_2$ ($-2$); only row $s_1$ has a positive entry (ratio
        $1$); pivot there:
        \begin{center}
        \begin{tabular}{c|rrrr|r}
        \hline
        Basis & $x_1$ & $x_2$ & $s_1$ & $s_2$ & RHS\\
        \hline
        $x_2$ & $-2$ & $1$ & $1$ & $0$ & $1$\\
        $s_2$ & $-1$ & $0$ & $1$ & $1$ & $3$\\
        \hline
        $z$ & $-5$ & $0$ & $2$ & $0$ & $2$\\
        \hline
        \end{tabular}
        \end{center}
        Now $x_1$ must enter (reduced cost $-5$), but its column holds
        $-2$ and $-1$: no positive entry, so the LP is \emph{unbounded}.
        Ray: $x_1=t\ge 0$ gives $x_2=1+2t$, $s_2=3+t$, i.e.
        \[
          \begin{pmatrix}x_1\\ x_2\end{pmatrix}
          =\begin{pmatrix}0\\ 1\end{pmatrix}
           +t\begin{pmatrix}1\\ 2\end{pmatrix},
          \qquad z=2+5t\;\longrightarrow\;\infty .
        \]
        (Check: $-2t+(1+2t)=1$ and $t-(1+2t)=-1-t\le 2$ for all
        $t\ge 0$.)

  \item The origin violates $2x_1+x_2\ge 8$: surplus $s_1$ and artificial
        $a$ in row~1, slack $s_2$ in row~2; phase one minimises $w=a$,
        i.e.\ $-w=-8+2x_1+x_2-s_1$:
        \begin{center}
        \begin{tabular}{c|rrrrr|r}
        \hline
        Basis & $x_1$ & $x_2$ & $s_1$ & $s_2$ & $a$ & RHS\\
        \hline
        $a$   & $2$ & $1$ & $-1$ & $0$ & $1$ & $8$\\
        $s_2$ & $1$ & $1$ & $0$ & $1$ & $0$ & $3$\\
        \hline
        $-w$ & $-2$ & $-1$ & $1$ & $0$ & $0$ & $-8$\\
        \hline
        \end{tabular}
        \end{center}
        Enter $x_1$ ($-2$); ratios $8/2=4$, $3/1=3$; $s_2$ leaves:
        \begin{center}
        \begin{tabular}{c|rrrrr|r}
        \hline
        Basis & $x_1$ & $x_2$ & $s_1$ & $s_2$ & $a$ & RHS\\
        \hline
        $a$   & $0$ & $-1$ & $-1$ & $-2$ & $1$ & $2$\\
        $x_1$ & $1$ & $1$ & $0$ & $1$ & $0$ & $3$\\
        \hline
        $-w$ & $0$ & $1$ & $1$ & $2$ & $0$ & $-2$\\
        \hline
        \end{tabular}
        \end{center}
        All entries $\ge 0$: phase one ends with $\min w=2>0$ and the
        artificial stuck at $a=2$: the LP is \emph{infeasible}.
        Interpretation of $w^{*}=2$: the largest value of $2x_1+x_2$
        compatible with $x_1+x_2\le 3$, $x\ge 0$ is $6$ (at $(3,0)$), so
        the requirement $2x_1+x_2\ge 8$ is missed by exactly $2$.
        Graphically the two constraint sets have no common point.

  \item Corners: $(0,0)$, $(3,0)$, $(2,1)$ (from $x_1+2x_2=4$ and
        $x_1+x_2=3$), $(0,2)$, with values of $z=x_1+c\,x_2$:
        \begin{center}
        \begin{tabular}{c|cccc}
        \hline
        corner & $(0,0)$ & $(3,0)$ & $(2,1)$ & $(0,2)$\\
        \hline
        $z$ & $0$ & $3$ & $2+c$ & $2c$\\
        \hline
        \end{tabular}
        \end{center}
        \begin{itemize}[leftmargin=1.6em,itemsep=0.3ex]
          \item $c=1$: $3=2+c$, tie between $(3,0)$ and $(2,1)$ with
                $z^{*}=3$ (while $(0,2)$ gives $2$); the objective is
                parallel to the edge $x_1+x_2=3$. Optimal set:
                $(x_1,x_2)=\lambda(3,0)+(1-\lambda)(2,1)
                =(2+\lambda,\,1-\lambda)$, $\lambda\in[0,1]$.
          \item $c=2$: $2+c=2c=4$, tie between $(2,1)$ and $(0,2)$ with
                $z^{*}=4$ (while $(3,0)$ gives $3$); parallel to
                $x_1+2x_2=4$. Optimal set:
                $(x_1,x_2)=\lambda(2,1)+(1-\lambda)(0,2)
                =(2\lambda,\,2-\lambda)$, $\lambda\in[0,1]$.
        \end{itemize}
        For $0\le c<1$ the unique optimum is $(3,0)$, for $1<c<2$ it is
        $(2,1)$, for $c>2$ it is $(0,2)$. Multiple optima exactly for
        $c\in\{1,2\}$.
\end{solutions}
